% INTERNAL DRAFT v45 for Scientific Reports. Generated by build_v40.py from team-approved JSONs. DO NOT SUBMIT without author approval.
\documentclass[11pt]{article}

\usepackage{mathptmx}
\usepackage[margin=1in]{geometry}
\usepackage{booktabs,tabularx,url,cite,amsmath,graphicx,lineno,setspace}
\linenumbers\onehalfspacing
\title{Gains from encrypted pooled recalibration at eICU hospitals with few calibration deaths depended on base-model class weighting and comparator}
\author{Shadman Mahmood Khan Pathan and Sachin Shetty\\
Department of Electrical and Computer Engineering, Old Dominion University, Norfolk, VA, USA\\
Correspondence: spath004@odu.edu}
\date{}
\begin{document}
\maketitle
\begin{abstract}
\noindent Risk models moved to a new hospital often need recalibration, but hospitals with few deaths cannot estimate even a simple local correction reliably. Using @@stays@@ intensive care stays from @@H@@ hospitals in eICU, we recalibrated a text-based mortality model at @@ntgt@@ target hospitals held out from predictor training but contributing calibration labels, and compared a 12-bin isotonic map fitted to bin summaries summed across hospitals under CKKS homomorphic encryption with a regularized local intercept. In the @@nlow@@ hospitals with 0--5 calibration deaths, the pooled map had lower mean Brier score (difference @@prim_mean@@ $\times10^{-3}$, 95\% interval @@prim_lo@@ to @@prim_hi@@) and was better at @@plower@@ hospitals. With an unweighted base model (post hoc), the difference was @@s_ub_loc_mean@@ $\times10^{-3}$ (@@s_ub_lo@@ to @@s_ub_hi@@). No incremental benefit over source-only recalibration was demonstrated. In 160 simulated health systems, a source-only map came close to the pooled map when targets resembled sources, and pooling harmed a growing share of low-event hospitals under heterogeneous prevalence. In a synthetic separate-process prototype with custodian-held keys, encryption took @@pr_enc@@~s and @@pr_mb@@~MB per hospital and maps matched plaintext. Pooled calibration should be benchmarked against source-based recalibration and under alternative base-model fitting before deployment.
\end{abstract}

\section*{Introduction}
A clinical prediction model has two jobs: to rank patients by risk and to state an absolute probability that clinicians can act on. Good ranking does not imply that predicted probabilities match observed event rates, and calibration often deteriorates when a model moves between hospitals with different case mix, documentation habits or outcome prevalence\cite{r3,vancalster2019,finlayson2021}. Recalibration with local outcome data is the standard remedy\cite{r17,janssen2008}. It is least reliable exactly where it is most needed: at small hospitals, or at hospitals where the outcome is rare, a local calibration set may contain only a handful of deaths.

Pooling calibration evidence across hospitals can stabilize the estimate, but it raises two concerns. The first is privacy: hospitals may be unwilling to share patient-level predictions and outcomes. Federated approaches keep patient records at each hospital and exchange only model updates or statistics\cite{rieke2020}. Distributed and federated calibration methods exchange aggregate statistics or model parameters\cite{r4,r5,r16,r6,r20}, and homomorphic encryption allows such statistics to be summed so that, with independent key custody and a controlled release policy, the aggregator need not see any single hospital's contribution\cite{r18,r7,r8}. Our eICU analysis demonstrates the arithmetic in a single process, and a separate-process prototype on synthetic hospitals tests key separation and cost (Methods). The second concern is statistical: a pooled map describes the pooled population and can move predictions away from the correct risk at a hospital whose prevalence differs from the rest.

Existing work has shown that distributed calibration is feasible and can improve average calibration\cite{r4,r5,r6}. Less is known about the hospitals for which pooling is proposed as most useful, those with very few observed events, and about whether any benefit comes from pooling itself or simply from using a more flexible calibration map. Here we address both questions with a prespecified, event-stratified comparison in the eICU Collaborative Research Database\cite{r1,r2}. Our primary analysis compares an encrypted pooled isotonic map with a regularized local intercept at target hospitals with 0--5 calibration deaths. Named controls separate pooling from map flexibility and from each hospital's own labels. Because true risks are unobservable in real data, we complement the eICU analysis with a protocol-fixed simulation that reruns the unmodified analysis code on synthetic health systems with known risks, and we map when pooling helps and when it harms. Post hoc, we also examine how the comparison depends on the class weighting of the base model.

\section*{Results}
\subsection*{Cohort and hospital partition}
The audited cohort comprised @@stays@@ selected intensive care unit stays from @@H@@ hospitals, with @@mort@@\% hospital mortality (@@deaths@@ deaths). Of @@landmark@@ selected unit stays with a unit-discharge offset of at least 24 hours, before text and hospital-size filters, the final cohort retained those at hospitals with at least 120 eligible records and sufficient recorded text; the stay-selection audit found @@promoted@@ stays promoted by missing fields. In the designated primary partition (seed 7), @@nsrc@@ hospitals trained the base model and @@ntgt@@ were targets for recalibration; @@nlow@@ target hospitals had 0--5 deaths in their calibration sets. Table~\ref{tab:eicu} summarizes the cohort and all prespecified contrasts. In the final cohort, recorded hospital discharge preceded the 24-hour landmark in @@d24@@ stays and preceded unit discharge in @@dunit@@; these recording inconsistencies were retained as specified and removed in a post hoc sensitivity analysis (below). Selected stays at source and target hospitals had age medians of @@age_src@@ and @@age_tgt@@ years, female proportions of @@fem_src@@\% and @@fem_tgt@@\%, and crude mortality of @@mort_src@@\% and @@mort_tgt@@\%; unit mix and 24-hour text length differed (median @@words_src@@ and @@words_tgt@@ words; Table~\ref{tab:splits}a).

\subsection*{Pooled recalibration at low-event hospitals}
Across the @@nlow@@ low-event target hospitals, the encrypted pooled map had a mean held-out Brier score of @@bpool4@@, compared with @@bloc4@@ for the regularized local intercept, about @@prim_rel@@\% lower for pooling (Fig.~\ref{fig:contrasts}a); with an unweighted base model the difference was small (see Post hoc sensitivity analyses). Across the held-out test stays of these hospitals, the local intercept predicted @@strat_pred_loc@@\% mean risk versus @@strat_obs@@\% observed mortality. The paired difference was @@prim@@ $\times10^{-3}$ (hospital bootstrap 95\% interval), with median @@prim_median@@ and interquartile range @@iqr_lo@@ to @@iqr_hi@@ $\times10^{-3}$. Pooling gave the lower Brier score at @@plower@@ hospitals and the local intercept at @@llower@@; no single-hospital deletion reversed the sign of the mean (@@delflip@@ of @@nlow@@). Excluding each hospital's own summary from the pooled map changed the observed contrast little (@@selfinc@@ $\times10^{-3}$). The advantage over the local intercept held for every ridge penalty from 0.5 to 8, ranging from @@pen_lo@@ to @@pen_hi@@ $\times10^{-3}$, and it grew with stronger shrinkage (Fig.~\ref{fig:contrasts}b).

\subsection*{Pooling versus source-based recalibration}
The prespecified controls qualify this result. Against a binned isotonic map fitted only to source hospitals, the contrast was @@srcbin@@ $\times10^{-3}$, and against source Platt scaling followed by a local intercept it was @@srcplatt@@ $\times10^{-3}$; both intervals included zero (Fig.~\ref{fig:contrasts}a). The leave-one-target-hospital-out pooled map gave the same pattern (@@lotoplatt@@ against source Platt plus local intercept). Pooled temperature scaling was clearly worse than the pooled binned map (@@pooltemp@@). Thus, at low-event hospitals, the encrypted pooled map outperformed a strongly shrunken local intercept, but no incremental benefit over recalibration that used source hospitals alone was demonstrated. These inconclusive intervals do not establish equivalence, because no equivalence margin was prespecified.

\subsection*{Sensitivity to the source--target split}
In each of the five split seeds fixed in the frozen analysis notebook, the pooled map had a lower mean Brier score than the regularized local intercept at target hospitals with 0--5 calibration deaths (@@spl_nmin@@ to @@spl_nmax@@ hospitals per split), with mean differences from @@spl_min@@ to @@spl_max@@ $\times10^{-3}$ (Table~\ref{tab:splits}b). Each unadjusted, within-split 95\% interval lay above zero; the overlapping splits are descriptive sensitivity analyses, not independent confirmatory tests. Over all @@ntgt@@ target hospitals the corresponding differences ranged from @@spl_all_min@@ to @@spl_all_max@@ $\times10^{-3}$. In every split, hospital-macro Brier scores of the pooled map and the source-only binned and Platt maps over all target hospitals differed by at most @@spl_macro_gap@@ $\times10^{-3}$ (Supplementary Table S2); these descriptive scores do not establish equivalence to source-based methods. The splits reuse patients and hospitals, and this two-group extract was prepared after the primary analysis.

\subsection*{Post hoc sensitivity analyses}
Two sensitivity analyses were specified after the frozen analysis and are post hoc (Table~\ref{tab:sens}). First, of the stays with inconsistent discharge offsets, @@s_order_prim@@ were at the @@nlow@@ primary hospitals and @@s_order_test@@ of these were in test sets. Removing flagged stays from evaluation only gave a difference of @@s_a1@@ $\times10^{-3}$ (paired change from the frozen analysis @@s_a1chg@@). Removing them before refitting the predictor and all recalibration maps, with each remaining stay's original split assignment kept, gave @@s_a2@@ $\times10^{-3}$ at the same @@nlow@@ hospitals (paired change @@s_a2chg@@). The positive direction therefore persisted at the original hospitals. After this exclusion a few additional hospitals had 0--5 calibration deaths; they were not added to the comparison.

Second, we refitted the base model without class weighting and refitted every recalibration map. In the primary-stratum test stays (@@s_ub_tstays@@ stays, @@s_ub_obs@@\% observed mortality), mean predicted risk was @@s_ub_unc@@\% before recalibration, @@s_ub_ppool@@\% after the pooled map and @@s_ub_ploc@@\% after the local intercept, and the four methods' mean Brier scores ranged from @@s_ub_bmin@@ to @@s_ub_bmax@@. The pooled-versus-local difference fell to @@s_ub_loc@@ $\times10^{-3}$ (paired change from the class-weighted model @@s_ub_chg@@ $\times10^{-3}$), and neither source-based contrast excluded zero. The pooled map's mean Brier score was similar under the two base models (@@s_bw_pool@@ and @@s_bu_pool@@), whereas the local intercept's fell from @@s_bw_loc@@ to @@s_bu_loc@@. The advantage over the local intercept was therefore highly sensitive to class weighting; numerically, most of the attenuation came from improved local-intercept performance after the predictor and calibrators were refitted. Because this refit changes more than the probability level, it does not isolate over-prediction as the cause, and similar Brier scores do not establish equivalence.

\subsection*{Calibration and discrimination}
Calibration diagnostics are consistent with the Brier difference. Over all @@ntgt@@ target hospitals (@@tstays@@ held-out stays, @@tdeaths@@ deaths), the pooled map had calibration intercept @@cal_pool_i@@ and slope @@cal_pool_s@@, close to source Platt scaling (@@cal_platt_i@@ and @@cal_platt_s@@), whereas the penalty-2 local intercept over-predicted, with intercept @@cal_loc_i@@ and slope @@cal_loc_s@@ (Fig.~\ref{fig:reliability}; an all-target diagnostic). The same pattern held in the low-event stratum (@@strat_stays@@ test stays, @@strat_deaths@@ deaths, @@strat_obs@@\% observed mortality): mean predicted risk was @@strat_pred_pool@@\% for the pooled map, @@strat_pred_src@@\% for the source-only binned map and @@strat_pred_loc@@\% for the local intercept, with calibration-in-the-large @@citl_pool@@, @@citl_src@@ and @@citl_loc@@ respectively. One interpretation, not tested directly, is that with few events the penalty keeps the local intercept close to the over-predicting class-weighted, source-temperature-scaled model. Discrimination was not the target of this comparison, but the binned isotonic map creates ties and can change it: across the five splits, hospital-macro AUROC over all target hospitals was @@auc_pool_min@@ to @@auc_pool_max@@ for the pooled map and @@auc_src_min@@ to @@auc_src_max@@ for the source-only binned map (@@auc_nlo@@ to @@auc_nhi@@ hospitals with both outcomes; Supplementary Table S2).

\subsection*{Encrypted aggregation}
Encrypted aggregation reproduced plaintext pooling. Rounded counts and deaths matched exactly, the temperature choice matched plaintext in all @@ckks_runs@@ splits (maximum loss-curve error @@ckks_loss@@), and encrypted and plaintext maps differed by at most @@ckks_map@@ over @@ckks_n@@ held-out predictions.

\subsection*{Separate-process prototype}
With each role in its own process and the secret key held only by a custodian, the aggregator's decryption attempt failed in all @@pr_runs@@ runs (3 per size, @@pr_hmin@@ to @@pr_hmax@@ synthetic hospitals, one @@pr_cpus@@-core machine), showing that it did not hold the key in this prototype, not security against malicious behavior. Rounded counts and deaths were exact, the temperature choice matched plaintext, and pooled maps differed from plaintext by at most @@pr_mapdiff@@. Per hospital, encrypting and serializing the five vectors took a median of @@pr_enc@@~s and produced @@pr_mb@@~MB of ciphertext, excluding the one-time @@pr_ctxmb@@~MB public context. Median aggregation times were @@pr_agg10@@, @@pr_agg50@@, @@pr_agg139@@ and @@pr_agg250@@~s for 10, 50, 139 and 250 hospitals, and decryption took about @@pr_dec@@~s (Supplementary Table S3); times exclude inter-process transfer. Packing all @@pr_packn@@ values into one ciphertext reduced each upload to @@pr_packmb@@~MB.

\subsection*{Simulation with known risks}\label{simulation-results}

Table~\ref{tab:sim} and Fig.~\ref{fig:sim} summarize all 160 replicates; three replicates had no hospital in the 0--5-death stratum and contribute no stratum contrast. With homogeneous prevalence ($\sigma=0$), pooling lowered the stratum-mean risk error relative to the local intercept in all 39 replicates with stratum hospitals (mean $8.7\times10^{-3}$), and only 1\% of low-event hospitals were harmed. At $\sigma=0.7$ the stratum mean still favored pooling in 97\% of replicates, but 22\% of low-event hospitals were harmed and the advantage over all target hospitals vanished. At $\sigma=1.0$ the stratum mean was near zero with a wide replicate range, 44\% of low-event hospitals were harmed, and pooling was worse than the local intercept across all target hospitals in every replicate. The Brier contrast pointed the same way but was noisier: at $\sigma=0$ it favored pooling in 82\% of replicates, whereas the risk-error contrast did so in all of them.

Two controls qualify the gain. Excluding the evaluated hospital's own summary changed the scenario-mean stratum risk error by at most $0.3\times10^{-3}$, so self-inclusion explained little of it. A source-only binned isotonic map came within $0.4$--$0.9\times10^{-3}$ of the pooled map, so most of the advantage over the local intercept came from the flexible binned map correcting the class-weighted probability scale, not from pooling across hospitals. The post hoc shift scenario separates the two. When target hospitals had lower mortality than sources ($\delta=-0.7$; about ten stratum hospitals per replicate), pooling beat the source-only map by $5.9\times10^{-3}$ [2.5, 10.7] and the local intercept by $14.6\times10^{-3}$ [6.0, 21.6] in every replicate, and 4\% of stratum hospitals were harmed. The source-only map was miscalibrated for the shifted targets. In an exploratory paired check within the same 40 replicates, the leave-one-target-hospital-out pooled map also beat the source-only map, by $5.9\times10^{-3}$ [2.5, 10.6] in every replicate, and self-inclusion contributed only $0.04\times10^{-3}$. The correction therefore came from other target hospitals' summaries, not from each hospital's own labels. When targets had higher mortality ($\delta=+0.7$), 21 of 40 replicates had no stratum hospital. In the 19 eligible replicates only, the pooled advantage over the source-only map was small and inconsistent ($1.6\times10^{-3}$, positive in 11 of 19) and 28\% of stratum hospitals were harmed. Across all runs, rounded CKKS counts and deaths matched plaintext exactly, the encrypted and plaintext temperature choices always agreed, the maximum loss-curve error was $1.0\times10^{-11}$, and the encrypted and plaintext maps differed by at most $3.8\times10^{-10}$ over 1,484,326 held-out predictions.



\section*{Discussion}
Hospitals with very few observed deaths are where local recalibration is least reliable and where pooling is most often proposed. Whether encrypted pooling helped them depended on the comparator and on the base model. Against a strongly regularized local intercept, the prespecified comparison clearly favored the encrypted pooled isotonic map at such hospitals, with the pooled map better at most hospitals and the result stable to deleting any one hospital and to the choice of ridge penalty. In the primary low-event hospitals, the class-weighted pipeline's local-intercept predictions averaged @@strat_pred_loc@@\% risk against @@strat_obs@@\% observed mortality, while pooled and source-only binned predictions averaged @@strat_pred_pool@@\% and @@strat_pred_src@@\%. This is consistent with a calibration problem for the strongly regularized local comparator, but does not isolate its cause. A post hoc refit with an unweighted base model brought the local intercept within $1\times10^{-3}$ of the pooled map while the pooled map's accuracy was essentially unchanged; that refit changed both the predictor and the calibrators. The same controls, however, did not demonstrate that this correction required pooling across hospitals: the comparators were a binned map fitted to source hospitals alone and source Platt scaling with a local intercept, and neither base-model analysis established an incremental benefit of pooling over these source-based comparators; the intervals included zero and, because no equivalence margin was prespecified, do not establish equivalence.

The simulation illustrates a possible mechanism for this pattern. When target hospitals resembled the source population, a source-only map came close to the pooled map, consistent with the eICU result. Among the simulated scenarios, pooling clearly beat a source-only map in the post hoc scenario in which target hospitals had lower mortality than sources, a situation that produces low-event hospitals and one in which the source population is the wrong reference. As prevalence became heterogeneous, a growing share of low-event hospitals was harmed by pooling even while the average still favored it. Mean improvements therefore cannot guarantee better risk estimates at a particular hospital, and hospital-level win/loss counts should be reported alongside averages.

These findings suggest a practical check for a health system that wants to recalibrate a model at low-event hospitals: a flexible map fitted to well-populated source hospitals is a simple baseline, and any benefit of encrypted pooling should be demonstrated against that baseline rather than only against a heavily shrunken local update. Established model-updating guidance likewise selects among simple and more extensive updates according to the data available\cite{janssen2008,vergouwe2017}. The size of any gain also depends on how the base model was fitted: in eICU, the advantage over the local intercept largely disappeared with an unweighted base model. In our simulation, pooling added value over the source-only map in the tested scenario where targets had lower mortality than sources; whether this holds in other health systems is untested. Encryption itself was not the bottleneck: CKKS aggregation of 12-bin summaries was numerically indistinguishable from plaintext, and in a separate-process prototype it cost about @@pr_enc@@~s and @@pr_mb@@~MB per hospital. Secret-sharing secure aggregation solves the same summation problem without homomorphic encryption and tolerates hospitals that drop out during a round\cite{bonawitz2017}; homomorphic encryption instead lets a single aggregator add contributions without pairwise exchanges between hospitals, at the cost of relying on an independent key holder. We did not compare the two, and both release the same pooled summaries, so the inference risks described in the threat model apply to either. Adding differential privacy noise to the released aggregates would require a specific mechanism and an assessment of its effect on calibration at low-event hospitals, which we leave to future work.

This study has limitations. It uses one retrospective database, one prespecified source--target split for the primary analysis and one text-based model built from structured diagnosis and treatment strings, so external validation in another health system is needed. The base model was developed for this study rather than an established severity score such as APACHE IVa, and whether the findings hold when recalibrating such a score was not tested. The landmark excludes early unit discharges, hospital-discharge status is not a time-limited endpoint, and a small fraction of stays had inconsistent discharge offsets; removing them in post hoc analyses left the direction of the primary contrast unchanged at the original hospitals. The primary advantage over the local intercept depended strongly on class weighting of the base model, and both sensitivity analyses were post hoc. The hospital bootstrap treats the fitted pipeline as fixed and does not add within-hospital sampling uncertainty. The eICU encryption ran all roles in one process; the separate-process prototype used synthetic hospitals on one machine and did not test network deployment, collusion, a malicious key holder or differencing attacks, so the study establishes arithmetic fidelity and key separation in a prototype, not deployed security or differential privacy. The pooled map is patient-weighted, and a hospital-equal variant was not evaluated. We did not assess performance across sociodemographic groups, so fairness and known health inequalities remain unexamined; the model uses only diagnosis and treatment text. Finally, no minimum important Brier difference was prespecified, and the analysis makes no claim about clinical utility.

In conclusion, encrypted pooled recalibration improved probabilistic accuracy over a regularized local intercept at hospitals with few observed deaths when the base model was class-weighted, with arithmetic that matched plaintext; with an unweighted base model the difference was small and its interval included zero. No incremental benefit over source-based recalibration was demonstrated in eICU. In the simulation, pooling beat a source-only map in a post hoc scenario in which target hospitals had lower mortality than source hospitals; we do not infer that this condition explains the eICU result.

\section*{Methods}
\subsection*{Study design and data source}\label{study-design-and-data-source}

This is a retrospective, multi-hospital prediction-model evaluation using eICU Collaborative Research Database version 2.0, which contains intensive care unit admissions in 2014 and 2015 at hospitals across the United States \cite{r1,r2}. The target is hospital-discharge mortality among patients whose selected eICU unit stay reaches the 24-hour landmark. The frozen loader first requires patient, hospital, unit-visit-number, and unit-discharge-offset fields, then sorts remaining rows by unique patient identifier, unit visit number, and stay identifier and retains one observed unit stay per patient. This pre-selection missing-field exclusion can select a different stay when an earlier ordered row lacks a required field; an aggregate audit that replays the selection rule counts such promoted stays (Table~\ref{tab:eicu}). Unit visit numbers identify order within a hospital admission, while eICU provides no systematic chronological ordering of a patient's separate hospital stays within the same discharge year \cite{r11}. We therefore call this a deterministic selected observed stay, not the first lifetime, hospital, or database ICU stay. Hospital-discharge status is subsequently required, the selected unit-discharge offset must be at least 1,440 minutes, and hospitals must have at least 120 eligible records with more than 20 characters of extracted text. The frozen code flags a nonmissing discharge status containing \texttt{expired}, \texttt{death}, or \texttt{dead} as positive; its agreement with documented status categories requires an aggregate code-frequency audit \cite{r11}; an aggregate diagnostic confirmed that every final-cohort status is \texttt{alive} or \texttt{expired} and that labels agree with \texttt{expired}.

Only diagnosis and treatment entries with offsets from zero through 1,440 minutes are used. These offsets record when an item was entered in eCareManager, so the restriction defines a recorded-entry-time landmark and does not establish clinical onset or fully prospective availability \cite{r13,r14}. The admissionDx field is omitted from the primary analysis because the eICU admission-diagnosis documentation permits entry during the first 36 hours and specifies an entered-time offset \cite{r12}. A missing diagnosis or treatment offset column raises an error, and entries without a recorded offset are excluded rather than admitted unwindowed. An empty treatment string cannot establish that a treatment was not administered \cite{r14}. Text features are built locally; the study does not upload eICU records to a hosted language-model service.

\subsection*{Model development and hospital partition}\label{model-development-and-hospital-partition}

The earlier project notebook specified five hospital-shuffling seeds (7, 11, 17, 23, 31), reused for this corrected analysis, with $\max(2,\lfloor 0.34H\rfloor)$ of $H$ eligible hospitals assigned as source hospitals and the remainder as targets. Within each hospital, about 50\% of eligible selected stays are reserved for model development, 25\% for calibration, and the remainder for untouched testing, using the seed-specific random partition. Only source-hospital development rows train the base classifier; target-hospital development rows are unused. Calibration rows and their outcome labels from both source and target hospitals contribute to the pooled map. The target hospitals are therefore held out from predictor training, not from calibration. This design evaluates collaborative recalibration after local outcomes have become available; it does not evaluate zero-label transfer to a new hospital. No test rows or test outcomes are used to train the predictor or calibrators. Seed 7 is the designated primary source--target split; the other seeds assess sensitivity and are not independent clinical replications. This analysis was specified before inspection of the corrected run, but was not prospectively registered and followed exploratory work on an earlier cohort.

The model uses a word and bigram TF-IDF vocabulary with sublinear term frequency (at most 30,000 features, minimum document frequency two) fitted on source development text, followed by L2-penalized logistic regression (liblinear, $C=1$) with balanced class weights. The uncalibrated decision score is converted to a probability by the logistic link. Class weighting can shift probability levels, so improvements relative to raw scores alone are not evidence of an advantage from pooling. No clinical language model is used.

\subsection*{Recalibration methods}\label{recalibration-methods}

For source-only temperature scaling, one positive temperature is fitted by bounded minimization of negative log likelihood ($0.05\le T\le 10$) on source-hospital calibration rows; the transformed prediction is sigmoid(z/T) \cite{r15}, where z is the base-model logit. Pooled temperature scaling sums each hospital\textquotesingle s loss at 96 prespecified, equally spaced candidate temperatures from 0.25 to 5.0 and selects the minimizing aggregate loss. The main pooled method divides uncalibrated probabilities into 12 fixed equal-width bins. Each hospital computes, for every bin b, its count n\_b, death count d\_b, and probability sum s\_b. These quantities are summed across hospitals; a weighted isotonic regression then maps the pooled bin means s\_b/n\_b to pooled event fractions d\_b/n\_b with weights n\_b. This fixed-bin weighted-isotonic calculation is an aggregate-statistic approximation: it need not equal individual-prediction isotonic regression and is not a reproduction of the distributed smooth-isotonic estimator in \cite{r4}. Predictions exactly equal to zero belong to the first bin.

The local comparators use the same hospital-specific calibration split. The code also fits per-hospital temperature and unbinned isotonic maps as exploratory comparators; single-class calibration sets may yield degenerate fits. The code also defines a local-only fallback (Supplementary Methods); because it can disadvantage hospitals with zero events, the primary comparator instead updates the source-temperature model by a hospital-specific intercept a, an established model-updating form \cite{r17,janssen2008}, minimizing the hospital\textquotesingle s sum of Bernoulli negative log likelihood plus $2a^2$ over -8 \textless= a \textless= 8. The fixed ridge coefficient is an implementation choice rather than a universal standard; a sensitivity analysis evaluates coefficients 0.5, 1, 2, 4, and 8 without selecting from target-test outcomes. Use of every available local calibration label is primary.

The prespecified analysis includes a source-only logistic calibration with a positive slope and intercept, followed by a local offset using the same local outcome labels. It also specifies a leave-one-target-hospital-out pooled isotonic map that excludes the evaluated hospital\textquotesingle s own calibration summary while retaining contributions from other hospitals. This control is computed on the single host by subtracting the target hospital's plaintext bin summary from the decrypted aggregate; it is an analytic control, not an encrypted protocol, and releasing both inclusive and leave-one-out aggregates in a deployment would expose that hospital's summary by differencing. To separate cross-hospital pooling from map flexibility, it adds source-only binned isotonic and pooled-temperature controls and a fixed local-intercept penalty grid. We report the designated primary contrast even when a control is discordant: a positive primary interval supports only a qualified comparison with the fixed local intercept, while a broader claim across named comparators requires the corresponding controls to support it. The self-including-versus-leave-one-out contrast is a diagnostic because no non-inferiority margin was prespecified.

\subsection*{Encryption scope and aggregation audit}\label{encryption-scope-and-aggregation-audit}

The revised notebook encrypts each hospital\textquotesingle s aggregate vector with a TenSEAL CKKS public context, homomorphically adds the ciphertexts, and requests decryption of the sum using a distinct private context \cite{r7,r18}. The study executes these roles in one process on one authorized data host. This is a cryptographic arithmetic emulation, not a deployment across independent institutions, an implemented trusted third party, or a formal differential privacy guarantee. A proposed deployment would require an independently controlled key custodian and protection against collusion, malicious reporting, and small-group differencing. The analysis ran with TenSEAL CKKS in the target environment; plaintext fallback is disabled.

The frozen source specifies CKKS polynomial modulus degree 8,192, coefficient moduli of 60/40/40/60 bits, and scale $2^{40}$. Per source--target split and backend, it encrypts five short vectors per participating hospital: 96 temperature-loss terms, one calibration count, and three 12-bin vectors for counts, deaths, and predicted-score sums. For $H$ hospitals the aggregation performs $5(H-1)$ ciphertext-vector additions and decrypts five pooled vectors. Runtime and message size were measured in the prototype below.

The eICU analysis ran all roles in one process, with plaintext vectors and the decryption key accessible to that process, so it cannot establish confidentiality against its host; decryption also reveals pooled bin-level summaries to the key holder, and small bins or overlapping releases could reveal more than intended (Supplementary Methods). Key separation and cost were therefore tested in a separate-process prototype on synthetic hospitals (below).

CKKS arithmetic is approximate. The frozen source checks that decrypted bin and event counts round within tolerance and rejects impossible counts; rounded zeros define empty bins. A supplementary seed-7 driver compares the count, normalized loss curve, bin statistics, and held-out map predictions against plaintext on the authorized host. Results are in Table~\ref{tab:eicu}. Arithmetic fidelity alone does not establish safe aggregate release for every hospital group size.

\subsection*{Threat model}\label{threat-model}
The intended deployment has three roles: participating hospitals, an aggregator that receives and adds ciphertexts, and an independent key custodian that holds the only CKKS secret key. We assume honest-but-curious parties that follow the protocol and no collusion between the aggregator and the custodian. Before release, the aggregator observes ciphertexts, message sizes and which hospitals took part; the custodian can decrypt the pooled vectors (pooled bin counts, deaths and predicted-score sums, and the pooled temperature-loss curve), from which the pooled map is computed; and the released map is visible to its recipients. Authenticated aggregation, a minimum number of participating hospitals, decryption restricted to approved pooled sums and independently governed key custody are deployment requirements that this study did not implement. The design does not protect against collusion, against hospitals that submit false summaries, or against inference from the pooled summaries themselves, for example by differencing releases computed over overlapping sets of hospitals or from bins containing very few patients, and it provides no differential privacy guarantee. The eICU analysis emulated all roles in one process and therefore establishes arithmetic fidelity only. The separate-process prototype tests key separation and numerical fidelity on one machine, but does not establish independently governed key custody or protection against selective decryption, collusion or differencing.

\subsection*{Separate-process prototype}\label{prototype-methods}

To test key separation and cost, we implemented the same aggregation with each role in a separate operating-system process that exchanged only serialized bytes, using synthetic hospitals and the frozen CKKS parameters. A key custodian created the context and the only secret key and published a public context without it. Hospital processes received only the public context and encrypted their five vectors; an aggregator added the ciphertexts, attempted decryption as a check, and forwarded the sums to the custodian, which decrypted them. We ran 3 replicates each with 10, 50, 139 and 250 hospitals and recorded encryption, aggregation and decryption time, ciphertext bytes, and agreement with a plaintext computation of the aggregates and pooled map (a test harness only; no party in a deployment would hold all plaintext vectors). The prototype ran on one machine (@@pr_cpus@@ CPU cores; Python @@pr_py@@, TenSEAL @@pr_ts@@); it does not measure network latency and does not address collusion, malicious inputs or differencing attacks.

\subsection*{Outcomes and analysis plan}\label{outcomes-and-analysis-plan}

The designated primary estimand is hospital-macro overall probabilistic accuracy, measured by the mean across target hospitals with 0--5 deaths in their full local calibration sets of $\Delta$Brier = Brier(regularized local intercept) - Brier(pooled binned isotonic) on the hospital\textquotesingle s untouched test set. Positive values favor the pooled binned-isotonic method against this named comparator. We report the number of contributing target hospitals, both absolute mean Brier scores, the mean and median within-hospital difference, interquartile range, the numbers of hospitals with strictly lower pooled, strictly lower local, and exactly equal Brier scores, and percentile endpoints from 2,000 paired hospital resamples for seed 7. No formal sample-size calculation was performed; the analysis includes all hospitals meeting its fixed eligibility and partition rules and imposes no five-hospital minimum. No minimum important Brier difference was prespecified; exclusion of zero would establish statistical direction for the named contrast, not clinical materiality. We do not convert the interval into an unadjusted $p$-value decision. As an exploratory robustness check computed from the frozen hospital pairs, we report how many single-hospital deletions reverse the sign of the mean difference; the leave-one-out means themselves are withheld because, with few hospitals, they reconstruct hospital-level differences. In the primary stratum many test sets may contain few or no deaths; for a test set with no deaths the Brier score equals the mean squared predicted risk, so differences there largely reflect predicted risk levels. Hospital-level scores remain on the authorized host. Repeated hospital appearances across alternate splits are correlated and are not independent observations. A host-side check reconciled each hospital's primary pooled and penalty-2 local scores with the scarcity and penalty tables.

The locked code computes test-set negative log likelihood, 12-bin equal-width expected calibration error (ECE), 12-bin equal-mass adaptive ECE, and area under the receiver operating characteristic curve. A separate pinned-source diagnostic computes descriptive calibration intercepts and slopes, exact-boundary prediction rates, and held-out reliability bins. These diagnostics pool all seed-7 target hospitals rather than the 0--5-death stratum, so they describe the maps overall and cannot by themselves establish calibration in the primary stratum. Brier score combines calibration and discrimination, so a lower Brier score alone does not establish a calibration-only effect \cite{r3}. A calibration-specific interpretation requires the diagnostic evidence; ECE is bin dependent and unstable in small event strata. Macro averages treat hospitals equally, while pooled-patient summaries weight large hospitals more and cannot replace the estimand. For the primary stratum itself, an exploratory check reports calibration-in-the-large (intercept with slope fixed at 1)\cite{vancalster2016} for the pooled, local-intercept, and source-only binned maps. The pooled map is record-weighted, so large calibration hospitals exert more influence even though the primary effect is hospital-macro. Released aggregates followed small-cell and complementary-suppression rules (Supplementary Methods); these are not a formal privacy guarantee. Primary inference is fixed before the audited test outputs; exploratory diagnostics receive no multiplicity-adjusted significance claim. Age above 89 years is top-coded in eICU and counted as 90 years or more.

The same patients may contribute to different hospital partitions across the five seeds, so split-level estimates are correlated. Site-level confidence intervals summarize between-hospital variation within the designated split; they do not account for the choice of source hospitals or establish external validity at future hospitals. Recalibration changes probability scale, while strictly monotone transformations preserve ranking and isotonic maps may create ties. No clinical risk threshold, treatment recommendation, or safety guarantee is evaluated in the primary analysis.

\textit{Post hoc sensitivity analyses.} After the frozen analysis, we specified two sensitivity analyses in a hash-pinned script that first reproduced every primary hospital's pooled and local Brier scores from the frozen run (tolerance $10^{-9}$) and the discharge-offset counts. (i) Stays whose recorded hospital discharge preceded the 24-hour landmark or unit discharge were removed either from test sets only or before refitting the predictor and all recalibration maps, with each remaining stay's original seed-7 assignment kept; contrasts were evaluated at the original @@nlow@@ primary hospitals. (ii) The base logistic model was refitted without class weighting, all other settings unchanged, and every recalibration method was refitted. Intervals use 2,000 paired hospital resamples with fixed seeds; paired changes compare each hospital's contrast with its frozen value. Counts of flagged stays from 1 to 19 were banded before release.

\subsection*{Supplementary simulation with known risks}\label{simulation-methods}

Because true risks are unobservable in eICU, we fixed a simulation protocol before running it and recorded its checksum. The simulation executes the hash-verified analysis code unmodified, apart from an in-memory return of fitted predictions: the TF-IDF model, 50/25/25 within-hospital split, 0.34 source fraction, split seed 7, 12-bin pooled isotonic map, penalty-2 local intercept, and TenSEAL CKKS aggregation are those of the locked analysis. Each replicate has 60 hospitals with sizes drawn log-normally around 450 (range 130--3,000) and hospital intercepts $a_h=\mathrm{logit}(0.09)+\sigma z_h$, $z_h\sim N(0,1)$. A patient with severity $x\sim N(0,1)$ has true risk $\mathrm{sigmoid}(a_h+1.5x)$. The text contains synthetic diagnosis tokens whose presence depends on $x$ and on a hospital-specific documentation shift, which creates cross-site drift. We ran 40 replicates at each $\sigma\in\{0,0.35,0.7,1.0\}$. The primary simulation contrast, averaged over target hospitals with 0--5 calibration deaths, is $\Delta$risk = MSE(local intercept, true risk) $-$ MSE(pooled map, true risk) on test rows; positive values favor pooling. We also report the Brier form of the eICU estimand, the leave-one-target-hospital-out and source-only binned-isotonic controls, and the share of low-event hospitals harmed by pooling. Replicate percentiles describe variation across synthetic worlds, not sampling uncertainty for any real hospital. After seeing these results, we added one post hoc scenario, also fixed and checksummed before it ran: with $\sigma=0.35$, every target hospital's intercept was shifted by $\delta\in\{-0.7,+0.7\}$ relative to source hospitals, with target membership taken in advance from the frozen seed-7 shuffle. It asks whether pooling adds anything beyond a source-only map when targets differ systematically from sources.



\subsection*{Reporting and ethics}
This retrospective study used only de-identified data from the eICU Collaborative Research Database, which is publicly available to credentialed researchers through PhysioNet\cite{r1}. The first author completed the human-research training required by PhysioNet and obtained credentialed access under its data use agreement before any analysis. All tables in the database are de-identified to meet the safe harbor provision of the US Health Insurance Portability and Accountability Act, and the re-identification risk was certified by an independent privacy firm (Privacert, Cambridge, MA; HIPAA certification no. 1031219-2)\cite{r2}. Because the study was a secondary analysis of publicly available, de-identified data with no contact with patients, review by the Old Dominion University institutional review board was not sought, and individual informed consent was not obtained, as the patients cannot be identified. All methods were carried out in accordance with relevant guidelines and regulations, including the PhysioNet data use agreement. Only aggregate results were exported from the authorized data host. Reporting follows TRIPOD+AI, with TRIPOD-Cluster considerations for the hospital structure\cite{r9,r19}; the completed TRIPOD+AI checklist is Supplementary Table S1. The analysis plan was fixed in a signed analysis lock verified before the run, and the simulation protocol was checksummed before it ran; the study was not prospectively registered. Large language model assistants (Claude, Anthropic; ChatGPT, OpenAI) were used for code review, auditing and manuscript editing under the authors' direction; the authors verified all analyses and take full responsibility for the content.

\section*{Data availability}
The eICU Collaborative Research Database v2.0 is available to credentialed users who complete the required training and sign the PhysioNet data use agreement (https://physionet.org/content/eicu-crd/2.0/); the authors cannot share record-level data. All eICU results in this article are aggregate and are reported in the tables, figures and Supplementary Table S2; prototype results are in Supplementary Table S3.

\section*{Code availability}
The data-independent code that produces all eICU results (the frozen analysis notebook, audit and diagnostic scripts, with dependency versions and checksums) the simulation code and the separate-process prototype will be deposited in a public repository before publication, with instructions for credentialed users, and is available to editors and reviewers during peer review.

\begin{thebibliography}{26}
\bibitem{r3} Huang Y, Li W, Macheret F, Gabriel RA, Ohno-Machado L. A tutorial on calibration measurements and calibration models for clinical prediction models. Journal of the American Medical Informatics Association. 2020;27(4):621--633. doi:10.1093/jamia/ocz228.
\bibitem{vancalster2019} Van Calster B, McLernon DJ, van Smeden M, Wynants L, Steyerberg EW; Topic Group `Evaluating diagnostic tests and prediction models' of the STRATOS initiative. Calibration: the Achilles heel of predictive analytics. BMC Medicine. 2019;17:230. doi:10.1186/s12916-019-1466-7.
\bibitem{finlayson2021} Finlayson SG, Subbaswamy A, Singh K, Bowers J, Kupke A, Zittrain J, Kohane IS, Saria S. The clinician and dataset shift in artificial intelligence. New England Journal of Medicine. 2021;385(3):283--286. doi:10.1056/NEJMc2104626.
\bibitem{r17} Steyerberg EW, Borsboom GJJM, van Houwelingen HC, Eijkemans MJC, Habbema JDF. Validation and updating of predictive logistic regression models: a study on sample size and shrinkage. Statistics in Medicine. 2004;23(16):2567-2586. doi:10.1002/sim.1844.
\bibitem{janssen2008} Janssen KJM, Moons KGM, Kalkman CJ, Grobbee DE, Vergouwe Y. Updating methods improved the performance of a clinical prediction model in new patients. Journal of Clinical Epidemiology. 2008;61(1):76--86. doi:10.1016/j.jclinepi.2007.04.018.
\bibitem{rieke2020} Rieke N, Hancox J, Li W, Milletar{\`i} F, Roth HR, et al. The future of digital health with federated learning. npj Digital Medicine. 2020;3:119. doi:10.1038/s41746-020-00323-1.
\bibitem{r4} Huang Y, Jiang X, Gabriel RA, Ohno-Machado L. Calibrating predictive model estimates in a distributed network of patient data. Journal of Biomedical Informatics. 2021;117:103758. doi:10.1016/j.jbi.2021.103758.
\bibitem{r5} Peng H, Yu H, Tang X, Li X. FedCal: Achieving Local and Global Calibration in Federated Learning via Aggregated Parameterized Scaler. Proceedings of the 41st International Conference on Machine Learning. PMLR. 2024;235:40331--40346. \url{https://proceedings.mlr.press/v235/peng24g.html}
\bibitem{r16} Cormode G, Markov IL. Federated Calibration and Evaluation of Binary Classifiers. Proceedings of the VLDB Endowment. 2023;16(11):3253-3265. doi:10.14778/3611479.3611523.
\bibitem{r6} Maddock S, Cormode G, Maple C. Private Federated Multiclass Post-hoc Calibration. arXiv preprint arXiv:2510.01987. 2025. \url{https://arxiv.org/abs/2510.01987}
\bibitem{r20} {\"O}ney B. EquiFL-X: Communication-constrained federated calibration via label-aggregated histograms for multi-view screening mammography under acquisition shift. Comput Methods Programs Biomed. 2026;286:109596. doi:10.1016/j.cmpb.2026.109596.
\bibitem{r18} Cheon JH, Kim A, Kim M, Song Y. Homomorphic Encryption for Arithmetic of Approximate Numbers. Advances in Cryptology - ASIACRYPT 2017. Lecture Notes in Computer Science 10624. Springer; 2017:409-437. doi:10.1007/978-3-319-70694-8\_15.
\bibitem{r7} Benaissa A, Retiat B, Cebere B, Belfedhal AE. TenSEAL: A Library for Encrypted Tensor Operations Using Homomorphic Encryption. arXiv preprint arXiv:2104.03152. 2021. \url{https://arxiv.org/abs/2104.03152}
\bibitem{r8} Pathan SMK, Wang Q, Takeshita J, Imran SB, Shetty S. Privacy-Preserving Post Deployment Model Calibration with Fully Homomorphic Encryption. 2026 International Conference on Computing, Networking and Communications (ICNC). IEEE; 2026:1--6. doi:10.1109/ICNC68183.2026.11416991.
\bibitem{r1} Pollard T, Johnson A, Raffa J, Celi LA, Badawi O, Mark R. eICU Collaborative Research Database (version 2.0). PhysioNet. 2019. doi:10.13026/C2WM1R. \url{https://physionet.org/content/eicu-crd/2.0/}
\bibitem{r2} Pollard TJ, Johnson AEW, Raffa JD, Celi LA, Mark RG, Badawi O. The eICU Collaborative Research Database, a freely available multi-center database for critical care research. Scientific Data. 2018;5:180178. doi:10.1038/sdata.2018.178.
\bibitem{vergouwe2017} Vergouwe Y, Nieboer D, Oostenbrink R, Debray TPA, Murray GD, Kattan MW, Koffijberg H, Moons KGM, Steyerberg EW. A closed testing procedure to select an appropriate method for updating prediction models. Statistics in Medicine. 2017;36(28):4529--4539. doi:10.1002/sim.7179.
\bibitem{bonawitz2017} Bonawitz K, Ivanov V, Kreuter B, Marcedone A, McMahan HB, Patel S, Ramage D, Segal A, Seth K. Practical secure aggregation for privacy-preserving machine learning. In: Proceedings of the 2017 ACM SIGSAC Conference on Computer and Communications Security. ACM; 2017:1175--1191. doi:10.1145/3133956.3133982.
\bibitem{r11} MIT Laboratory for Computational Physiology. eICU Collaborative Research Database: patient table documentation. \url{https://eicu.mit.edu/eicutables/patient/} (accessed 24 September 2026).
\bibitem{r13} MIT Laboratory for Computational Physiology. eICU Collaborative Research Database: diagnosis table documentation. \url{https://eicu.mit.edu/eicutables/diagnosis/} (accessed 24 September 2026).
\bibitem{r14} MIT Laboratory for Computational Physiology. eICU Collaborative Research Database: treatment table documentation. \url{https://eicu.mit.edu/eicutables/treatment/} (accessed 24 September 2026).
\bibitem{r12} MIT Laboratory for Computational Physiology. eICU Collaborative Research Database: admissionDx table documentation. \url{https://eicu.mit.edu/eicutables/admissiondx/} (accessed 24 September 2026).
\bibitem{r15} Guo C, Pleiss G, Sun Y, Weinberger KQ. On Calibration of Modern Neural Networks. Proceedings of the 34th International Conference on Machine Learning. PMLR. 2017;70:1321-1330. \url{https://proceedings.mlr.press/v70/guo17a.html}
\bibitem{vancalster2016} Van Calster B, Nieboer D, Vergouwe Y, De Cock B, Pencina MJ, Steyerberg EW. A calibration hierarchy for risk models was defined: from utopia to empirical data. Journal of Clinical Epidemiology. 2016;74:167--176. doi:10.1016/j.jclinepi.2015.12.005.
\bibitem{r9} Collins GS, Moons KGM, Dhiman P, Riley RD, Beam AL, Van Calster B, et al. TRIPOD+AI statement: updated guidance for reporting clinical prediction models that use regression or machine learning methods. BMJ. 2024;385:e078378. doi:10.1136/bmj-2023-078378.
\bibitem{r19} Debray TPA, Collins GS, Riley RD, Snell KIE, Van Calster B, Reitsma JB, Moons KGM. Transparent reporting of multivariable prediction models developed or validated using clustered data: TRIPOD-Cluster checklist. BMJ. 2023;380:e071018. doi:10.1136/bmj-2022-071018.
\end{thebibliography}

\section*{Funding}
This research received no specific grant from any funding agency in the public, commercial or not-for-profit sectors.

\section*{Patient and public involvement}
Patients and the public were not involved in the design, conduct or reporting of this study.

\section*{Author contributions}
S.M.K.P. designed the study, wrote the code, ran the analyses and drafted the manuscript. S.S. supervised the study, reviewed the analysis plan and results, and revised the manuscript. Both authors approved the final version.

\section*{Competing interests}
The authors declare no competing interests.

\clearpage
\begin{table}[p]
\caption{eICU cohort and seed-7 results. $\Delta$Brier is the named comparator minus the first-named method, $\times10^{-3}$, with paired hospital-bootstrap 95\% interval (2,000 resamples); positive favors the first-named method.}
\label{tab:eicu}
\centering\footnotesize
\begin{tabularx}{\textwidth}{@{}p{.33\textwidth}p{.19\textwidth}X@{}}
\toprule
Measure & Result & Detail \\
\midrule
@@EicuTableRows@@
\bottomrule
\end{tabularx}
\end{table}

\begin{table}[p]
\caption{Synthetic simulation, 40 replicates per $\sigma$ and per post hoc $\delta$. Contrasts are local minus pooled ($\times10^{-3}$); positive favors pooling.}
\label{tab:sim}
\centering\footnotesize
\setlength{\tabcolsep}{2.2pt}
\begin{tabular}{@{}lccccccc@{}}
\toprule
$\sigma$ & Reps & Hosp. & $\Delta$risk [2.5, 97.5\%] & $>0$ & $\Delta$Brier ($>0$) & Harmed & All \\
\midrule
0.00 & 39 & 3.3 & 8.7 [2.7, 17.1] & 100\% & 8.0 (82\%) & 1\% & 6.4 \\
0.35 & 39 & 3.6 & 8.2 [1.4, 15.3] & 97\% & 9.3 (90\%) & 11\% & 4.6 \\
0.70 & 39 & 4.7 & 6.8 [0.5, 15.0] & 97\% & 8.5 (97\%) & 22\% & $-0.3$ \\
1.00 & 40 & 6.0 & 0.9 [$-16.3$, 19.1] & 60\% & 0.7 (55\%) & 44\% & $-7.6$ \\
\midrule
$-0.7^\dagger$ & 40 & 9.8 & 14.6 [6.0, 21.6] & 100\% & 14.0 (100\%) & 4\% & 9.0 \\
$+0.7^\dagger$ & 19 & 0.7 & 1.9 [$-3.7$, 7.1] & 79\% & 3.2 (53\%) & 28\% & $-0.2$ \\
\bottomrule
\end{tabular}

\vspace{2pt}
{\scriptsize Reps: replicates with at least one 0--5-death target hospital; every contrast column, including All, averages over these replicates. Hosp.: mean number of 0--5-death target hospitals over all 40 replicates. $>0$: share of those replicates favoring pooling. Harmed: mean share of 0--5-death hospitals where pooling increased risk error. All: $\Delta$risk over all target hospitals. $^\dagger$Post hoc: $\sigma=0.35$ with target-hospital shift $\delta$.}
\end{table}

\begin{table}[p]
\caption{Cohort characteristics and split sensitivity. \textbf{a}, Final cohort by hospital role in the primary split (seed 7). \textbf{b}, Paired $\Delta$Brier, local intercept minus pooled map ($\times10^{-3}$; positive favors pooling), with paired hospital-bootstrap 95\% intervals, for each of the five fixed split seeds. Splits reuse patients and are correlated.}
\label{tab:splits}
\centering\footnotesize
\textbf{a}\quad\begin{tabular}[t]{@{}lrrr@{}}
\toprule
 & All & Source & Target \\
\midrule
@@CohortRows@@
\bottomrule
\end{tabular}

\vspace{2pt}
{\scriptsize @@CohortNote@@}

\vspace{8pt}
\textbf{b}\quad\begin{tabular}[t]{@{}lrrrr@{}}
\toprule
Seed & $n$ (0--5) & 0--5 calibration deaths & $n$ (all) & All target hospitals \\
\midrule
@@SplitRows@@
\bottomrule
\end{tabular}

\vspace{2pt}
{\scriptsize $^*$Designated primary split; its 0--5 interval is the locked primary interval. Other intervals use split-specific bootstrap seeds.}
\end{table}

\begin{table}[p]
\caption{Post hoc sensitivity analyses at the @@nlow@@ primary hospitals (0--5 calibration deaths, seed 7). Differences are comparator minus pooled Brier ($\times10^{-3}$; positive favors pooling) with paired hospital-bootstrap 95\% intervals; two decimals with signs are shown when an interval includes zero. \textbf{a}, Mean Brier score and contrasts with the class-weighted (primary) and unweighted base models; all maps refitted. \textbf{b}, Local intercept minus pooled map after removing stays with inconsistent discharge offsets.}
\label{tab:sens}
\centering\footnotesize
\textbf{a}\quad\begin{tabular}[t]{@{}lrrrr@{}}
\toprule
 & \multicolumn{2}{c}{Mean Brier} & \multicolumn{2}{c}{Comparator $-$ pooled} \\
Method & Weighted & Unweighted & Weighted & Unweighted \\
\midrule
@@SensRowsA@@
\bottomrule
\end{tabular}

\vspace{2pt}
{\scriptsize Paired change in local minus pooled, unweighted minus weighted: @@s_ub_chg@@ $\times10^{-3}$. Scores within $1\times10^{-3}$ do not establish equivalence.}

\vspace{8pt}
\textbf{b}\quad\begin{tabular}[t]{@{}lrr@{}}
\toprule
Analysis & Local $-$ pooled & Change from frozen \\
\midrule
@@SensRowsB@@
\bottomrule
\end{tabular}
\end{table}

\begin{figure}[p]
\centering
\includegraphics[width=\textwidth]{fig2_eicu_contrasts.pdf}
\caption{Paired Brier contrasts at the @@nlow@@ eICU target hospitals with 0--5 calibration deaths (seed 7). \textbf{a}, Mean difference (comparator minus first-named method) with hospital-bootstrap 95\% intervals; blue intervals exclude zero. \textbf{b}, Pooled map versus the local intercept across ridge penalties.}
\label{fig:contrasts}
\end{figure}

\begin{figure}[p]
\centering
\includegraphics[width=0.55\textwidth]{fig3_reliability.pdf}
\caption{Held-out reliability over all @@ntgt@@ eICU target hospitals (seed 7), pooling all target stays; this is an all-target diagnostic, not specific to the 0--5-death primary stratum. Twelve equal-width bins; bins with fewer than 20 stays are suppressed. The dotted line marks perfect calibration.}
\label{fig:reliability}
\end{figure}

\begin{figure}[p]
\centering
\includegraphics[width=0.8\textwidth]{fig_sim_boundary_v25.pdf}
\caption{Simulation boundary. Left: mean $\Delta$risk (local minus pooled) with 2.5--97.5\% replicate ranges, for low-event targets and all targets. Right: share of low-event hospitals harmed by pooling.}
\label{fig:sim}
\end{figure}
\end{document}
